\documentclass[11pt,a4paper]{article}
\usepackage[T1]{fontenc}
\usepackage{lmodern}
\usepackage[margin=27mm]{geometry}
\usepackage{amsmath,amssymb,amsthm,mathtools}
\usepackage{graphicx,booktabs,microtype}
\usepackage[numbers,sort&compress]{natbib}
\usepackage[colorlinks=true,linkcolor=blue,citecolor=blue,urlcolor=blue]{hyperref}
\hypersetup{pdftitle={Exact two-particle relaxation crossover for the symmetric inclusion process on hypercubes},
 pdfauthor={Xiaojian Zeng},pdfsubject={Two-particle reversible relaxation on finite hypercubes},
 pdfkeywords={symmetric inclusion process, spectral gap, hypercube, two-particle dynamics}}
\newtheorem{theorem}{Theorem}
\newtheorem{lemma}{Lemma}
\newcommand{\SIP}{\mathrm{SIP}}
\newcommand{\ind}{\mathbf{1}}
\newcommand{\cube}{Q_D}
\newcommand{\ip}[2]{\langle #1,#2\rangle}
\title{Exact two-particle relaxation crossover for the\protect\\
 symmetric inclusion process on hypercubes}
\author{Xiaojian Zeng\thanks{\href{mailto:xiaojianzeng@foxmail.com}{xiaojianzeng@foxmail.com}}\protect\\[2pt]
 \small PBC School of Finance, Tsinghua University\protect\\
 \small Beijing 100190, China}
\date{}
\begin{document}
\maketitle
\begin{abstract}
The symmetric inclusion process describes diffusion with hopping rates
that increase with the occupation of the destination. Its slowest relaxation mode
can differ from that of a single particle. We determine the full spectral
gap for exactly two particles on every finite Boolean hypercube, with a
homogeneous positive site parameter and equal positive edge conductances.
The gap is the minimum of the inherited one-particle eigenvalue and the
first positive eigenvalue of an explicit birth-death chain for the Hamming
distance between the particles. Translation characters exhaust the
two-particle state space: every nontrivial character is bounded below by
the one-particle gap, while all nonradial modes in the translation-invariant
sector lie higher. A rank-one resolvent calculation then gives an exact
integral formula for the finite-parameter crossover. A slower separation
mode occurs precisely below this positive threshold in dimensions four
and higher; the four-dimensional threshold is $1/35$. A recurrence determines
the threshold in every dimension and proves its inverse-square dimension
scaling. The result specifies when two-particle relaxation is controlled
by collective separation rather than a linear occupation observable.
It concerns a finite graph with exactly two particles and leaves the
general many-particle degree-two gap conjecture unresolved.
\end{abstract}
\noindent\textbf{Keywords:} symmetric inclusion process; spectral gap;
 hypercube; reversible relaxation; two-particle dynamics.

\section{Introduction}
The symmetric inclusion process is an established attractive interacting
particle model. For conservative dynamics, spectral inheritance makes
one-particle eigenvalues available in every particle sector, but does not
ensure that they govern relaxation. Kim and Sau proved the one-particle
gap identity when all site parameters are at least one
\cite{KimSau2024}, and subsequently established its failure on suitable
geometries outside that regime \cite{KimSau2026}. They also proved an
asymptotic two-particle reduction as diffusivity vanishes along a fixed
positive direction. Their proposed exact identity
$\lambda_k=\lambda_2$ for all $k\geq2$, arbitrary finite connected graphs
and positive site parameters remains a conjecture. The later work of
Caputo, Quattropani and Sau proves degree-two reduction for stochastic
exchange dynamics and explicitly retains the Brownian energy
process/inclusion-process question \cite[Sec.~1.4.6]{CaputoQuattropaniSau2026}.

Here we calculate the full two-particle gap on a coherent graph family
at finite positive diffusivity. For the Boolean hypercube $Q_D$, we obtain
an exact boundary between one-particle and pair-separation relaxation.
The contribution is the full-sector comparison and the finite-parameter
boundary in every dimension. Qualitative failure of one-particle reduction
is already known from \cite{KimSau2026}; we do not claim that phenomenon
as new. Relative coordinates for two interacting particles likewise have
prior applications, including the exact one-dimensional transition
calculation of Carinci, Giardin\`a and Redig \cite{CarinciGiardinaRedig2017}.
The present calculation concerns the finite hypercube and its full
two-particle spectral gap. It gives no identity at particle number three
or higher.

\section{Model and exact result}
Write $V=(\mathbb Z/2\mathbb Z)^D$, $D\geq1$, and join $x$ to $x+e_i$
for each coordinate vector $e_i$. Addition is modulo two. All edges have
conductance $c>0$, and the site parameter is $a>0$. For fixed particle
number $k$, the state space is
$\Xi_k=\{\eta\in\mathbb N_0^V:\sum_x\eta_x=k\}$ and the backward
generator is
\begin{equation}\label{eq:generator}
 (L_k f)(\eta)=c\sum_{x\in V}\sum_{i=1}^D
 \eta_x(a+\eta_{x+e_i})
 \bigl[f(\eta-\delta_x+\delta_{x+e_i})-f(\eta)\bigr].
\end{equation}
Terms with $\eta_x=0$ vanish. This directed jump sum has no factor
$1/2$. It conserves particle number, has no reservoirs, and is reversible
with probability law
\begin{equation}\label{eq:measure}
 \mu_{a,k}(\eta)=\frac{k!}{(2^D a)_k}
 \prod_{x\in V}\frac{(a)_{\eta_x}}{\eta_x!},
 \qquad (a)_m=\prod_{j=0}^{m-1}(a+j),\quad (a)_0=1.
\end{equation}
These conventions agree with the conservative generator and measure of
\cite[Eqs.~(2.1)--(2.2)]{KimSau2026}.
Let $\lambda_k$ be the smallest strictly positive eigenvalue of $-L_k$
in $L^2(\mu_{a,k})$. All parameters except $c$ are dimensionless;
$c$, the eigenvalues and the gaps have inverse-time units. The relaxation
time of the slowest nonconstant mode is $1/\lambda_k$.

For $d\in\{0,\ldots,D\}$, define
\begin{equation}\label{eq:rates}
 \ell_d=2da+2\,\ind_{\{d=1\}},\qquad u_d=2(D-d)a.
\end{equation}
The negative birth-death generator $T_D(a)$ has action
\begin{equation}\label{eq:radial}
 (T_D(a)h)_d=\ell_d(h_d-h_{d-1})+u_d(h_d-h_{d+1}),
\end{equation}
with absent endpoint terms omitted. Write $\rho_D(a)$ for its first
positive eigenvalue. Conductance has been set to one in this auxiliary
matrix, so that the corresponding physical eigenvalue is $c\rho_D(a)$.

\begin{theorem}\label{thm:main}
For every finite integer $D\geq1$ and every $a,c>0$, the model
\eqref{eq:generator} satisfies
\begin{equation}\label{eq:gap}
 \lambda_1=2ca,\qquad
 \lambda_2=c\min\{2a,\rho_D(a)\}.
\end{equation}
Define
\begin{equation}\label{eq:threshold}
 a_D=(D-1)2^{-(D-1)}\int_0^1(1+u^2)^{D-1}\,du-1.
\end{equation}
For $D\leq3$, $\lambda_2=\lambda_1$ at every $a>0$. For $D\geq4$,
$a_D>0$ and
\begin{equation}\label{eq:crossing}
 \lambda_2<\lambda_1\ \Longleftrightarrow\ 0<a<a_D,
 \qquad
 \lambda_2=\lambda_1\ \Longleftrightarrow\ a\geq a_D.
\end{equation}
The thresholds obey
\begin{equation}\label{eq:recurrence}
 a_{D+1}=\frac{D^2 a_D+1}{(D-1)(2D+1)}\quad(D\geq2),
 \qquad a_4=\frac1{35},\qquad \lim_{D\to\infty}D^2a_D=1.
\end{equation}
\end{theorem}

The proof first controls every symmetry sector of the full occupation
space, then locates the crossing inside the remaining distance chain.
This order is necessary: distance lumping alone supplies an upper bound
on $\lambda_2$, rather than the equality in \eqref{eq:gap}.

\section{Full two-particle sector comparison}
Set $c=1$ throughout the proof. A two-particle configuration is an unordered
pair $\{x,y\}$, including the stack $x=y$. The same generator is obtained
from labelled positions by giving the first particle a jump $x\to v$ at
rate $a+\ind_{\{y=v\}}$ on an edge, and similarly for the second particle.
On symmetric functions the labelled and occupation generators agree.
In particular, a stack has total rate $2a$ to each neighboring destination.

\subsection{Translation characters}
For $q\in V$, let $\chi_q(x)=(-1)^{q\cdot x}$. Translations preserve the
generator and reversible measure, so their character spaces are invariant,
orthogonal, and exhaustive. A function in character $q$ can be written as
\begin{equation}\label{eq:character}
 F(\{x,y\})=\chi_q(x)g(x+y).
\end{equation}
Exchanging $x$ and $y$ requires
$g(z)=\chi_q(z)g(z)$. Thus $g$ is supported on
$A_q=\{z:\chi_q(z)=1\}$. Conversely, this support restriction makes
\eqref{eq:character} well-defined. With $N=2^D$, the character dimensions
are $N$ at $q=0$ and $N/2$ at each $q\ne0$, and their sum is
\begin{equation}
 N+(N-1)\frac N2=\frac{N(N+1)}2=|\Xi_2|.
\end{equation}
There is therefore no missing two-particle sector.

Substitution of the labelled jump rates gives the negative-generator
block on $A_q$:
\begin{equation}\label{eq:block}
 (B_q g)(z)=\bigl(2Da+2\,\ind_{\{|z|=1\}}\bigr)g(z)
 -2\sum_{i:q_i=0}\bigl(a+\ind_{\{z=e_i\}}\bigr)g(z+e_i),
 \qquad z\in A_q,
\end{equation}
where $|z|$ is Hamming weight. If $q_i=1$, the two off-diagonal
contributions cancel, and their departure rates remain in the diagonal.
The extra rate at $z=e_i$ describes attraction into a stack. At $z=0$
the diagonal is $2Da$, with no extra attraction term.

\begin{lemma}\label{lem:nonzero}
For $q\ne0$, every eigenvalue of $B_q$ is at least $2|q|a$.
Characters with $|q|=1$ attain $2a$.
\end{lemma}
\begin{proof}
Let $m=|q|\geq1$ and $r=D-m$. The $m$ coordinates where $q_i=1$
have even parity in $A_q$ and are unchanged by surviving off-diagonal
moves. Fix their bit pattern. If all those bits are zero, the remaining
component is
\begin{equation}\label{eq:components}
 B_0(Q_r)+2ma I.
\end{equation}
Here $B_0(Q_r)$ is itself a translation-invariant block of a reversible
Markov negative generator, with nonnegative spectrum and a zero
eigenvalue. For $r=0$ it denotes the one-dimensional zero operator.
Any other allowed pattern has at least two nonzero bits, so that neither
$z=0$ nor $|z|=1$ occurs. Its component is
$2a\mathcal L_{Q_r}+2ma I$, where $\mathcal L_{Q_r}$ is the ordinary
cube Laplacian. This proves the lower bound. In the all-zero component
with $m=1$, the constant function in the remaining coordinates attains
$2a$.
\end{proof}

The attained eigenfunctions include the linear occupation observables
\begin{equation}\label{eq:linear}
 \Phi_i(\eta)=\sum_{x\in V}(-1)^{x_i}\eta_x.
\end{equation}
Their eigenvalue is $2a$. Independently, the one-particle negative
generator is $a\mathcal L_{Q_D}$; Walsh characters have Laplacian
eigenvalues $2|q|$, giving $\lambda_1=2a$.

\subsection{The translation-invariant sector}
The block $B_0$ is reversible for the unnormalized weights
\begin{equation}\label{eq:relative-weights}
 w(0)=\frac{a+1}{a},\qquad w(z)=1\quad(z\ne0).
\end{equation}
One can obtain these weights directly from \eqref{eq:measure}:
distinct-position configurations have unnormalized mass $a^2$ and stacks
have mass $a(a+1)/2$, while their translation orbits have sizes $N/2$
and $N$, respectively. Detailed balance on $0\leftrightarrow e_i$
also gives $w(0)2a=w(e_i)2(a+1)$.

Coordinate permutations preserve $w$ and $B_0$. Their invariant subspace
$\mathcal R$ consists of functions of $|z|$. Its weighted orthogonal
complement consists of functions with zero sum on every Hamming shell.
In particular, for $f\in\mathcal R^\perp$,
\begin{equation}\label{eq:nonradial-conditions}
 f(0)=0,\qquad \sum_{|z|=1}f(z)=0,\qquad \sum_z f(z)=0.
\end{equation}
Self-adjointness makes this complement invariant. Its weighted norm
equals its counting norm, and \eqref{eq:block} reduces to
\begin{equation}\label{eq:nonradial-action}
 B_0 f=2a\mathcal L_{Q_D}f+2\,\ind_{\{|z|=1\}}f.
\end{equation}
At zero, the equality uses the unit-shell sum in
\eqref{eq:nonradial-conditions}. The counting-space gap of
$\mathcal L_{Q_D}$ is two. Consequently,
\begin{equation}\label{eq:nonradial-bound}
 \frac{\ip{f}{B_0 f}_w}{\ip{f}{f}_w}
 =\frac{2a\ip{f}{\mathcal L_{Q_D}f}+2\sum_{|z|=1}|f(z)|^2}
 {\sum_z|f(z)|^2}\geq4a
 \quad(0\ne f\in\mathcal R^\perp).
\end{equation}

On $\mathcal R$, there are $d$ flips decreasing distance and $D-d$
increasing it, with the additional rate two into zero from distance one.
The restriction is therefore exactly \eqref{eq:radial}. Its reversible
masses may be chosen as
\begin{equation}\label{eq:radial-weights}
 \pi_0=\frac{a+1}{a},\qquad \pi_d=\binom Dd\quad(1\leq d\leq D).
\end{equation}
All neighboring transition rates are positive, so this finite chain is
irreducible. Lemma~\ref{lem:nonzero}, \eqref{eq:nonradial-bound}, and
the attained eigenvalue $2a$ prove the two-particle formula
\eqref{eq:gap} at $c=1$.

\section{Exact location of the crossover}
\subsection{A radial resolvent at the competing eigenvalue}
Let $R_D$ be the ordinary radial cube Laplacian: its diagonal is $D$
and its downward and upward entries are $-d$ and $-(D-d)$. Its eigenvalues
are $0,2,\ldots,2D$. For example, the shell sums of Walsh characters of
each weight give $D+1$ radial eigenfunctions with these distinct
eigenvalues. With $\varepsilon_d$ the coordinate vectors of
$\mathbb R^{D+1}$, \eqref{eq:rates} gives
\begin{equation}\label{eq:rankone}
 \frac{T_D(a)}{2a}=R_D+a^{-1}\varepsilon_1
 (\varepsilon_1-\varepsilon_0)^{\mathsf T}.
\end{equation}
Since $M=R_D-I$ is invertible, the determinant lemma yields
\begin{equation}\label{eq:detlemma}
 \det\bigl(T_D(a)-2aI\bigr)
 =(2a)^{D+1}\det(M)\left(1+\frac{h_D}{a}\right),\qquad
 h_D=(\varepsilon_1-\varepsilon_0)^{\mathsf T}M^{-1}\varepsilon_1.
\end{equation}
Put $v=M^{-1}\varepsilon_1$. The row at distance zero gives
$(D-1)v_0-Dv_1=0$, hence $h_D=-v_0/D$.
On the full cube, the unit-shell indicator is
$\operatorname{Adj}(Q_D)\delta_0=(DI-\mathcal L_{Q_D})\delta_0$.
The corresponding radial equation therefore implies
\begin{equation}\label{eq:green}
 v_0=(D-1)G_D-1,\qquad
 G_D=\bigl[(\mathcal L_{Q_D}-I)^{-1}\bigr]_{00}
 =2^{-D}\sum_{j=0}^D\frac{\binom Dj}{2j-1}.
\end{equation}
The last expression follows by normalized Walsh diagonalization;
none of its denominators is zero.

Writing the positive-$j$ terms as integrals gives
\begin{equation}
 G_D=2^{-D}\left[-1+
 \int_0^1\frac{(1+u^2)^D-1}{u^2}\,du\right].
\end{equation}
Integration by parts has a vanishing boundary term at $u=0$ and yields
\begin{equation}\label{eq:greenintegral}
 G_D=-1+D2^{-(D-1)}\int_0^1(1+u^2)^{D-1}\,du.
\end{equation}
Combining \eqref{eq:green}--\eqref{eq:greenintegral} proves $h_D=-a_D$.
Thus a radial eigenvalue can equal $2a$ at a positive parameter only when
$a=a_D$. If $a_D>0$, the zero of
\eqref{eq:detlemma} as a function of $a$ is simple.

\subsection{Which eigenvalue crosses, and in which direction?}
For every $a>0$, $T_D(a)$ is similar by \eqref{eq:radial-weights} to a
real symmetric tridiagonal matrix with nonzero neighboring entries. Its
eigenvalues are real, simple and continuous in $a$, and exactly one is
zero. Moreover, $T_D(a)/a\to2R_D$ as $a\to\infty$. Every positive radial
eigenvalue is therefore larger than $2a$ for sufficiently large $a$.

If $a_D\leq0$, the determinant has no positive-parameter zero, so none
of these eigenvalues can cross $2a$; in particular $\rho_D(a)>2a$
for all $a>0$. Suppose instead $a_D>0$. A higher positive radial
eigenvalue cannot be the one equal to $2a$ at $a_D$: the first positive
one would then be strictly below $2a$ there, forcing another crossing
between $a_D$ and the large-$a$ regime. This contradicts the unique
possible zero of \eqref{eq:detlemma}. Hence the crossing belongs to
$\rho_D$. The simple determinant zero, together with eigenvalue
simplicity, enforces a sign change. Continuity and the absence of any
other crossing show that
\begin{equation}\label{eq:rhorange}
 \rho_D(a)\begin{cases}
 <2a,&0<a<a_D,\\
 =2a,&a=a_D,\\
 >2a,&a>a_D.
 \end{cases}
\end{equation}
No monotonicity assumption on $\rho_D(a)/a$ is needed.

\subsection{Threshold signs, recurrence and dimension scaling}
Define $H_m=2^{-m}\int_0^1(1+u^2)^m\,du$. Integrating the derivative
of $u(1+u^2)^m$ gives
\begin{equation}\label{eq:Hrecurrence}
 H_m=\frac{1+mH_{m-1}}{2m+1},\qquad H_0=1.
\end{equation}
Since $a_D=(D-1)H_{D-1}-1$, this identity implies
\eqref{eq:recurrence}. Direct evaluation gives
\begin{equation}\label{eq:values}
 a_1=-1,\quad a_2=-\frac13,\quad a_3=-\frac1{15},\quad
 a_4=\frac1{35},\quad a_5=\frac{17}{315},\quad a_6=\frac{37}{693}.
\end{equation}
The recurrence propagates positivity from $D=4$. Together with
\eqref{eq:rhorange} and \eqref{eq:gap}, this proves the crossover
statements.

For the asymptotic, put $b_D=D^2a_D$. Then
\begin{equation}\label{eq:brecurrence}
 b_{D+1}=q_D(b_D+1),\qquad
 q_D=\frac{(D+1)^2}{(D-1)(2D+1)}\longrightarrow\frac12.
\end{equation}
For $D\geq6$, $0<q_D\leq4/5$, so the recurrence first bounds $b_D$
uniformly. Also,
\begin{equation}
 |b_{D+1}-1|\leq\frac45|b_D-1|+|2q_D-1|.
\end{equation}
Iteration gives a geometrically decaying initial contribution and a
convolution of a summable geometric sequence with a term tending to
zero. Both vanish as $D\to\infty$, proving $b_D\to1$.
Finally, restoring $c$ multiplies every eigenvalue by $c$ and leaves
$a_D$ unchanged. This completes the proof of Theorem~\ref{thm:main}.

\section{Four-dimensional example and reproducibility}
For $D=4$, the radial characteristic polynomial, with the convention
$p_4(\lambda)=\det(\lambda I-T_4(a))$, is
\begin{align}\label{eq:poly}
 p_4(\lambda)&=\lambda(\lambda-8a)
 \bigl[\lambda^3-(32a+2)\lambda^2
 +(304a^2+32a)\lambda\nonumber\\[-2pt]
 &\hspace{39mm}{}-768a^3-48a^2\bigr].
\end{align}
In particular $p_4(2a)=96a^4(35a-1)$. Below $a_4=1/35$,
$\rho_4(a)$ is the smallest positive root of the cubic factor, and
$\lambda_2=c\rho_4(a)<2ca$. Figure~\ref{fig:crossover} illustrates
this crossing and the dimension scaling of the exact threshold.

\begin{figure}[tbp]
 \centering
 \includegraphics[width=\textwidth]{manuscript/figures/crossover.pdf}
 \caption{The exact two-particle reduction. Left: the radial gap divided
 by the one-particle gap on $Q_4$, and the full two-particle ratio obtained
 by taking its minimum with one. The vertical line is the exact crossing
 $a/a_4=1$, with $a_4=1/35$. Eigenvalues are evaluated at 80-decimal working
 precision for illustration; the crossing is established analytically.
 Right: $D^2a_D$ from the rational recurrence, for $4\leq D\leq64$,
 with its proved limiting value one. Neither panel asserts a
 thermodynamic phase transition.}
 \label{fig:crossover}
\end{figure}

The accompanying symbolic verification compares full occupation
generators assembled from direct jump rates with a separate reversible
incidence construction. It checks character embeddings, radial
intertwiners and exact root ordering at $D=3$, $a=1/100$, and at $D=4$,
$a=1/100,1/35,1/10,1$. Resolvent, integral and recurrence identities are
also compared exactly for $1\leq D\leq16$. These finite checks support
the calculation; the analytic sector bounds and crossing argument
establish the unrestricted finite-$D$ statement. The figure-generation
script writes the plotted data and compares high-precision radial
eigenvalues with a separately assembled double-precision symmetric
matrix. Those numerical values are illustrations, not exact certificates.
Regeneration instructions and package versions are supplied with the
reproducibility material.

\section{Relaxation modes and scope}
Below $a_D$, a translation-invariant function of particle separation is
the slowest two-particle mode. The bounds on the other sectors ensure
that its slower relaxation is not an artifact of radial projection.
At and above the threshold, a linear occupation observable such as
\eqref{eq:linear} attains the gap. The result therefore distinguishes
pair-separation relaxation from a mode already visible to a single
particle, and locates the distinction at finite diffusivity.

Published spectral monotonicity gives
$\lambda_k\leq\lambda_2$ for every $k\geq2$
\cite[Eq.~(2.7)]{KimSau2026}. Hence below the threshold one-particle
reduction is excluded in all those sectors. This observation uses an
existing monotonicity result and gives no equality between $\lambda_k$
and $\lambda_2$. Above the threshold, the many-particle gap may still be
smaller than the one-particle gap. The general degree-two conjecture
is thus left unresolved.

All conclusions assume exactly two particles for the calculated gap,
strictly positive homogeneous site parameter, equal positive edge
conductances, and a finite closed cube. Heterogeneous weights, other
geometries and reservoir dynamics are outside the result. The limit
$D^2a_D\to1$ describes the sequence of two-particle crossover boundaries;
it is not a statement about a thermodynamic transition or a uniform
many-particle limit.

\section*{Data, code and use of generative AI}
Reproducibility scripts, exact verification outputs, plotted data and
environment requirements accompany this manuscript. No experimental
dataset is used. OpenAI Codex was used in the research process for
candidate exploration, analytic derivations, code generation, symbolic
and numerical checking, literature comparison, and manuscript drafting.
Separate AI-assisted reconstruction and checks were performed; they do
not constitute external peer review. AI tools are not authors. Human
scientific review and approval of this draft remain pending.

\begingroup
\small
\setlength{\bibsep}{4pt}
\bibliographystyle{unsrtnat}
\bibliography{manuscript/references}
\endgroup
\end{document}
